\quad A realização deste trabalho permitiu compreender o funcionamento dos
algoritmos \textit{Insertion Sort}, \textit{Bubble Sort}, \textit{Selection
Sort}, \textit{Shell Sort}, \textit{Merge Sort} e \textit{Quick Sort}, este
último em duas versões, além de analisar seus comportamentos em diferentes
tipos de entrada. Por meio das implementações, dos testes de mesa, das
análises de complexidade e dos experimentos práticos, foi possível observar
as diferenças entre as estratégias utilizadas por cada algoritmo.

\quad A análise teórica mostrou que o \textit{Insertion Sort} apresenta
complexidade $\Theta(n)$ no melhor caso, quando o vetor já está ordenado, e
$\Theta(n^2)$ no caso médio e no pior caso. O \textit{Bubble Sort} e o
\textit{Selection Sort}, considerando as implementações utilizadas neste
trabalho, apresentam complexidade $\Theta(n^2)$ nos três casos analisados.
Já o \textit{Shell Sort} possui comportamento dependente da sequência de
\textit{gaps} utilizada. Na implementação estudada, com a divisão sucessiva
do \textit{gap} por 2, o melhor caso apresenta crescimento de ordem
$\Theta(n\log n)$, enquanto o pior caso apresenta complexidade
$\Theta(n^2)$. O \textit{Merge Sort} apresenta complexidade temporal
$\Theta(n\log n)$ no melhor caso, no caso médio e no pior caso, além de
necessitar de memória auxiliar proporcional ao tamanho do vetor. Por fim, o
\textit{Quick Sort} apresenta complexidade $\Theta(n\log n)$ no melhor caso
e no caso médio, e $\Theta(n^2)$ no pior caso, que ocorre quando o pivô
gera partições muito desbalanceadas. Por ordenar o vetor no próprio local e
realizar a chamada recursiva apenas sobre a menor partição, a implementação
utilizada exige memória auxiliar de ordem $\Theta(\log n)$.

\quad Os experimentos realizados com entradas crescentes, decrescentes e
aleatórias permitiram observar na prática as diferenças previstas pela
análise teórica. No \textit{Insertion Sort}, a organização inicial dos
elementos possui influência significativa no tempo de execução: o tempo é
praticamente nulo quando o vetor já está ordenado e aproximadamente o dobro
do observado em entradas aleatórias quando ele está em ordem decrescente.
No \textit{Selection Sort}, o tempo foi semelhante nos três tipos de
entrada, pois a quantidade de comparações não depende da organização
inicial. O \textit{Bubble Sort} apresentou os maiores tempos em entradas
decrescentes e aleatórias, mas, devido à interrupção antecipada quando não
há trocas, teve tempo desprezível em entradas crescentes. O \textit{Shell
Sort} e o \textit{Merge Sort} apresentaram tempos muito inferiores aos dos
algoritmos quadráticos nos três tipos de entrada, o que evidencia a
vantagem dos algoritmos de custo próximo de $\Theta(n\log n)$ para
instâncias grandes.

\quad Os experimentos com o \textit{Quick Sort} evidenciaram a importância
da escolha do pivô. A versão 1, que utiliza o primeiro elemento, apresentou
tempos próximos aos dos algoritmos quadráticos em entradas crescentes e
decrescentes, pois nesses casos o pivô é sempre o menor ou o maior valor da
partição, levando o algoritmo ao pior caso. Em entradas aleatórias, porém,
seu tempo foi comparável ao dos algoritmos de custo $\Theta(n\log n)$. A
versão 2, que utiliza a média aritmética dos elementos como pivô, manteve
as partições equilibradas nos três tipos de entrada, com tempos muito
baixos, ao custo de um percurso adicional a cada particionamento. Vale
ressaltar que esse bom desempenho depende da distribuição dos valores
utilizados nos testes, pois a média aritmética pode gerar partições
desbalanceadas quando os valores são muito discrepantes.

\quad Também foi possível observar que o aumento do tamanho da entrada
provoca um crescimento no tempo de execução dos algoritmos. Esse
crescimento se torna mais significativo nos algoritmos e situações em que
há maior quantidade de comparações, trocas ou deslocamentos. Dessa forma,
os resultados experimentais possibilitam relacionar o comportamento
observado nos gráficos e tabelas com as complexidades obtidas na análise
teórica.

\quad No caso do \textit{Shell Sort}, teoricamente, a escolha do divisor
influencia a sequência de \textit{gaps} utilizada: um divisor menor produz
uma quantidade maior de etapas, enquanto um divisor maior reduz essa
quantidade. Entretanto, a quantidade de etapas, por si só, não determina o
tempo de execução, pois a sequência de \textit{gaps} também influencia a
quantidade de elementos que precisam ser comparados e deslocados. Como a
alteração do divisor não pôde ser implementada, os experimentos deste
trabalho utilizaram apenas o divisor 2, e a comparação entre diferentes
sequências de \textit{gaps} permanece como possibilidade de trabalho
futuro.

\quad Dessa forma, o trabalho permitiu relacionar conceitos teóricos de
análise de algoritmos com resultados obtidos por meio de experimentos
práticos. A comparação entre os seis algoritmos, considerando as duas
versões do \textit{Quick Sort}, mostrou que, apesar de todos possuírem o
mesmo objetivo de ordenar os elementos, suas estratégias de execução
produzem comportamentos diferentes de acordo com o tamanho e a organização
inicial da entrada. Os algoritmos de custo $\Theta(n\log n)$ se destacaram
em instâncias grandes, desde que, no caso do \textit{Quick Sort}, a escolha
do pivô evite partições desbalanceadas.